\documentclass[openany]{book}
\input{notes_white}
\usepackage{mathtools}
\usepackage{circuitikz}
\settitle{The Mysteries of Recursion}
\begin{document}
\title{\Large{\underemphasise{Basic Algebra}} \\ \vspace{5pt}
\large{Course Notes MA2205: Basic Algebra \\ \vspace{5pt} Notes by Prof. Anirban Banerjee}}
\author{Eklavya Raman \\ \vspace{5pt} 24MS024} 
\date{\today}
\maketitle
\tableofcontents
\chapter*{Introduction}
    Around the year 1800, people already knew how to solve equations of the forms - 
    \begin{itemize}
        \item Linear Equations: \( ax + b = 0 \)
        \item Quadratic Equations: \( ax^2 + bx + c = 0 \)
        \item Cubic Equations: \( ax^3 + bx^2 + cx + d = 0 \)
        \item Quartic Equations: \( ax^4 + bx^3 + cx^2 + dx + e = 0 \)
    \end{itemize}
    However, the general solution for equations of degree 5 or higher was not known. A young mathematician named Evariste Galois answered the question of the existence of a general solution for polynomial equations of degree 5 or higher, using a tool called groups. \\\\
    Around the same time, Carl Friedrich Gauss introduced a technique called modular arithmetic, which laid the foundation for number theory. This technique shares a close relationship with groups. \\\\
    The 1800s also saw a revolution in geometry. For over 2000 years, Euclidean Geometry was the only known geometry. However, mathematicians began to explore different geometries (specailly geometries on non-planar or curved surfaces). Groups became useful in studying these new geometries as well. \\\\
    The formalization of groups by mathematicians supporting Galois and Gauss led to the birth of Group Theory and Abstract Algebra. Mathematicians realized that groups were very useful in studying various mathematical structures, and thus extended the idea of groups to other structures such as rings, fields, modules, vector spaces etc. It took years of effort by many mathematicians to formalize these structures and study their properties. \\\\
    As time passed, the applications of these structures in various fields of mathematics and science became apparent. Today, it is used in computer science, cryptography, physics, chemistry, biology and many other fields. \\\\
    This course aims to introduce the basic concepts of Group Theory, Ring Theory and some fundamental ideas in Abstract Algebra. We will explore the properties of these structures and their applications in various fields of mathematics and science.

\chapter{Groups}
Before we formally define groups, we consider a few examples to get a basic idea of what groups are and how they arise in various mathematical contexts.
\section{Examples of Groups}
\begin{enumerate}
    \item \textbf{Clock Arithmetic:} Consider a planet which spins faster than the Earth, such that a day on this planet is 7 hours long. The time on this planet can be represented using numbers from 0 to 6, where 0 represents midnight, 1 represents 1 AM, and so on. Let us visualise this using a clock with 7 hours.
    \begin{center}
    \begin{circuitikz}[scale=1]
        \draw (0,0) circle(2cm);
        \foreach \i in {0,...,6} {
            \draw[thick] ({90 - \i * 360/7}:1.8cm) -- ({90 - \i * 360/7}:2cm);
            \node at ({90 - \i * 360/7}:1.5cm) {\i};
        }
        \draw[thick,->] (0,0) -- ({90 - 2 * 360/7}:1.3cm) node[midway, right] {};
    \end{circuitikz}
    \end{center}
    Let us define an operation of addition on this clock, where adding a number \( a \) to the current time \( b \) results in moving \( a \) hours forward on the clock. For example, if the current time is 5 and we add 4 hours, we move 4 hours forward on the clock, resulting in a time of 2 (since we wrap around after reaching 6). Let us visualise this operation on the clock.
    \begin{center}
    \begin{circuitikz}[scale=1]
        \draw (0,0) circle(2cm);
        \foreach \i in {0,...,6} {
            \draw[thick] ({90 - \i * 360/7}:1.8cm) -- ({90 - \i * 360/7}:2cm);
            \node at ({90 - \i * 360/7}:1.5cm) {\i};
        }
        \draw[thick,->] (0,0) -- ({90 - 5 * 360/7}:1.3cm) node[midway, right] {};
        \draw[thick,->,red] ({90 - 5 * 360/7}:0.5cm) arc[start angle={90 - 5 * 360/7}, end angle={90 - (5+4) * 360/7}, radius=0.5cm];
        \draw[thick,->] (0,0) -- ({90 - 2 * 360/7}:1.3cm) node[midway, right] {};
    \end{circuitikz}
    \end{center}
    Here, we see that $5 + 4 = 2$. Similarly, we can define subtraction on this clock, where subtracting a number \( a \) from the current time \( b \) results in moving \( a \) hours backward on the clock. For example, if the current time is 2 and we subtract 5 hours, we move 5 hours backward on the clock, resulting in a time of 4 (since we wrap around after reaching 0). Let us visualise this operation on the clock.
    \begin{center}
    \begin{circuitikz}[scale=1]
        \draw (0,0) circle(2cm);
        \foreach \i in {0,...,6} {
            \draw[thick] ({90 - \i * 360/7}:1.8cm) -- ({90 - \i * 360/7}:2cm);
            \node at ({90 - \i * 360/7}:1.5cm) {\i};
        }
        \draw[thick,->] (0,0) -- ({90 - 2 * 360/7}:1.3cm) node[midway, right] {};
        \draw[thick,->,red] ({90 - 2 * 360/7}:0.5cm) arc[start angle={90 - 2 * 360/7}, end angle={90 - (2-5) * 360/7}, radius=0.5cm];
        \draw[thick,->] (0,0) -- ({90 - 4 * 360/7}:1.3cm) node[midway, right] {};
    \end{circuitikz}
    \end{center}
    Notice that $2 - 5 = 4$, and also that subtraction can be thought of as adding the negative of a number, i.e., $2 - 5 = 2 + (-5) = 2 + 2 = 4$. \\\\
    This type of arithmetic is called modular arithmetic, specifically arithmetic modulo 7 in this case. The set of numbers \( \{0, 1, 2, 3, 4, 5, 6\} \) with the operations of addition and subtraction defined as above forms a mathematical structure known as a group. In this group, the operation of addition is closed, associative, has an identity element (0), and every element has an inverse (for example, the inverse of 5 is 2 since \( 5 + 2 \equiv 0 \mod 7 \)). \\\\
    Modular arithmetic has applications in various fields such as cryptography, computer science, and number theory.
    \item \textbf{Symmetries of an Equilateral Triangle:} Consider an equilateral triangle with vertices labeled \( A, B, C \). The symmetries of this triangle include rotations and reflections that map the triangle onto itself. Let us visualize these symmetries.
    \begin{center}
    \begin{circuitikz}[scale=1]
        \draw (0,0) -- (2,0) -- (1,1.732) -- cycle;
        \node at (0,-0.3) {C};
        \node at (2,-0.3) {B};
        \node at (1,2.0) {A};
    \end{circuitikz}
    \end{center}
    The operations that map the triangle onto itself are:
    \begin{itemize}
        \item \textbf{Rotations:} We can rotate the triangle about its center by \( 120^\circ \) clockwise, to get the triangle mapped onto itself. [Denoted by $R$.]
        \item \textbf{Flipping:} Flipping the triangle about an axis passing through a vertex and the midpoint of the opposite side results in mapping the triangle onto itself. [Denoted by $F$.]
        \item \textbf{No Operation:} The triangle remains unchanged [Denoted by $I$].
    \end{itemize}
    We can compose these fundamental operations to get all the symmetries of the triangle. Let us visualise these operations.
    \begin{center}
    \begin{circuitikz}[scale=1]
        \draw (0,0) -- (2,0) -- (1,1.732) -- cycle;
        \node at (0,-0.3) {C};
        \node at (2,-0.3) {B};
        \node at (1,2.0) {A};
        \draw[thick,->] (3,1) -- (4,1) node[midway, above] {$R$};
        \draw (5,0) -- (7,0) -- (6,1.732) -- cycle;
        \node at (5,-0.3) {B};
        \node at (7,-0.3) {A};
        \node at (6,2.0) {C};
        \draw[thick,->] (8,1) -- (9,1) node[midway, above] {$R$};
        \draw (10,0) -- (12,0) -- (11,1.732) -- cycle;
        \node at (10,-0.3) {A};
        \node at (12,-0.3) {C};
        \node at (11,2.0) {B};
    \end{circuitikz}
    \end{center}
    Similarly, we can visualize the flipping operation.
    \begin{center}
    \begin{circuitikz}[scale=1]
        \draw (0,0) -- (2,0) -- (1,1.732) -- cycle;
        \node at (0,-0.3) {C};
        \node at (2,-0.3) {B};
        \node at (1,2.0) {A};
        \draw[thick,->] (3,1) -- (4,1) node[midway, above] {$F$};
        \draw (5,0) -- (7,0) -- (6,1.732) -- cycle;
        \node at (5,-0.3) {B};
        \node at (7,-0.3) {C};
        \node at (6,2.0) {A};
    \end{circuitikz}
    \end{center}
    Let us now compose the flipping operation with rotations to get more symmetries.
    \begin{center}
    \begin{circuitikz}[scale=1]
        \draw (0,0) -- (2,0) -- (1,1.732) -- cycle;
        \node at (0,-0.3) {C};
        \node at (2,-0.3) {B};
        \node at (1,2.0) {A};
        \draw[thick,->] (3,1) -- (4,1) node[midway, above] {$R$};
        \draw (5,0) -- (7,0) -- (6,1.732) -- cycle;
        \node at (5,-0.3) {B};
        \node at (7,-0.3) {A};
        \node at (6,2.0) {C};
        \draw[thick,->] (8,1) -- (9,1) node[midway, above] {$F$};
        \draw (10,0) -- (12,0) -- (11,1.732) -- cycle;
        \node at (10,-0.3) {A};
        \node at (12,-0.3) {C};
        \node at (11,2.0) {B};
    \end{circuitikz}
    \end{center}
    Continuing this process, we find that there are a total of 6 symmetries of the equilateral triangle: \( I, R, R^2, F, FR, FR^2 \). Note that double flipping results in the identity operation, i.e., \( F^2 = I \) and that \( R^3 = I \). Let us visualise the above two operations to verify that they result in the identity operation.
    \begin{center}
    \begin{circuitikz}[scale=1]
        \draw (0,0) -- (2,0) -- (1,1.732) -- cycle;
        \node at (0,-0.3) {C};
        \node at (2,-0.3) {B};
        \node at (1,2.0) {A};
        \draw[thick,->] (3,1) -- (4,1) node[midway, above] {$F$};
        \draw (5,0) -- (7,0) -- (6,1.732) -- cycle;
        \node at (5,-0.3) {B};
        \node at (7,-0.3) {C};
        \node at (6,2.0) {A};
        \draw[thick,->] (8,1) -- (9,1) node[midway, above] {$F$};
        \draw (10,0) -- (12,0) -- (11,1.732) -- cycle;
        \node at (10,-0.3) {C};
        \node at (12,-0.3) {B};
        \node at (11,2.0) {A};
    \end{circuitikz}
    \end{center}
    \begin{center}
    \begin{circuitikz}[scale=0.8]
        \draw (0,0) -- (2,0) -- (1,1.732) -- cycle;
        \node at (0,-0.3) {C};
        \node at (2,-0.3) {B};
        \node at (1,2.0) {A};
        \draw[thick,->] (3,1) -- (4,1) node[midway, above] {$R$};
        \draw (5,0) -- (7,0) -- (6,1.732) -- cycle;
        \node at (5,-0.3) {B};
        \node at (7,-0.3) {A};
        \node at (6,2.0) {C};
        \draw[thick,->] (8,1) -- (9,1) node[midway, above] {$R$};
        \draw (10,0) -- (12,0) -- (11,1.732) -- cycle;
        \node at (10,-0.3) {A};
        \node at (12,-0.3) {C};
        \node at (11,2.0) {B};
        \draw[thick,->] (13,1) -- (14,1) node[midway, above] {$R$};
        \draw (15,0) -- (17,0) -- (16,1.732) -- cycle;
        \node at (15,-0.3) {C};
        \node at (17,-0.3) {B};
        \node at (16,2.0) {A};
    \end{circuitikz}
    \end{center}
\item \textbf{Integers under Addition:} Consider the set of integers \( \mathbb{Z} = \{\ldots, -3, -2, -1, 0, 1, 2, 3, \ldots\} \) with the operation of addition. \\
\textbf{Note: } We can define subtraction as adding the negative of a number, i.e., $x - y = x + (-y)$. We now have the following properties of this set equipped with the addition operation -
\begin{itemize}
    \item \underline{Closure:} Let $x, y \in \mathbb{Z}$, then $x + y \in \mathbb{Z}$.
    \item \underline{Additive Identity:} There is a unique number \( 0 \in \mathbb{Z} \) such that for any \( x \in \mathbb{Z} \), \( x + 0 = x \). This number is called the identity element.
\end{itemize}
\end{enumerate}
\begin{exercise}
    Instead of an equilateral triangle, consider the follwing shapes and find all their symmetries - 
    \begin{itemize}
        \item A square
        \item A rectangle (not a square)
        \item A cube
        \item A Tetrahedron
    \end{itemize}
\end{exercise}
\underline{Soln:} Let us consider just the case of a square here. \\ 
\textbf{Square:} Consider a square with vertices labeled \( A, B, C, D \). We can do the operations of rotation by \( 90^\circ \) clockwise (denoted by $R$), Flipping about the vertical axis ($F_v$), flipping about diagonal axis ($F_d$) and no operation ($I$). We can compose these operations to get all the symmetries of the square. There are a total of 8 symmetries of the square. \\
Similarly, we can find the symmetries of the other shapes. \\\\
Let us now analyze the similarities between these examples. 
\begin{enumerate}
    \item Each of these examples has a set of objects that are considered (numbers of the clock, symmetries of the triangle, integers). We call these objects elements of a set.
    \item We have a way of combining two elements of the set to get another element of the set (addition on the clock, composition of symmetries, addition of integers). We call this way of combining elements an operation.
    \item All sets are closed under the operation defined on them (adding two numbers on the clock gives another number on the clock, composing two symmetries gives another symmetry, adding two integers gives another integer).
    \item The operation is associative (for any three numbers on the clock, adding them in any order gives the same result, composing three symmetries in any order gives the same result, adding three integers in any order gives the same result).
    \item There is an identity element in each set with respect to the operation (0 on the clock, identity symmetry of the triangle, 0 in integers).
    \item Every element has an inverse with respect to the operation (for any number on the clock, there is another number that when added gives 0, for any symmetry of the triangle, there is another symmetry that when composed gives the identity symmetry, for any integer, there is another integer that when added gives 0).
\end{enumerate}
\section{Groups and Operations}
\subsection{Definitions}
Based on the observations from the previous section, we can now formally define groups and operations.
\begin{mydefinition}{Groups}
    A group \( (G, *) \) is a set \( G \) equipped with a \underemphasise{binary operation} \( * : G \times G \to G \) that satisfies the following properties:
    \begin{itemize}
        \item \textbf{Closure:} For all \( a, b \in G \), the result of the operation \( a * b \) is also in \( G \).
        \item \textbf{Associativity:} For all \( a, b, c \in G \), \( (a * b) * c = a * (b * c) \).
        \item \textbf{Identity Element:} There exists an element \( e \in G \) such that for every element \( a \in G \), \( e * a = a * e = a \).
        \item \textbf{Inverse Element:} For each \( a \in G \), there exists an element \( b \in G \) such that \( a * b = b * a = e \), where \( e \) is the identity element.
    \end{itemize}
\end{mydefinition}
It took years of effort to develop the formal definition of groups based on the most fundamental properties observed in various mathematical contexts. For example, the property of commutativity (i.e., \( a * b = b * a \) for all \( a, b \in G \)) is not included in the definition of groups, as there exist many important groups where this property does not hold. Groups where commutativity holds are called abelian groups, named after the mathematician Niels Henrik Abel. Similarly, a non-commutative group is called non-abelian. \\
Let us now formally define a binary operation. 
\begin{mydefinition}{Binary Operation}
    A binary operation on a set \( S \) is a function \( * : S \times S \to S \) that combines any two elements \( a, b \in S \) to produce another element \( c \in S \), denoted as \( a * b = c \). 
\end{mydefinition}
Let us consider some examples of binary operations.
\begin{enumerate}
    \item \textbf{Addition on Integers:} The operation of addition \( + : \mathbb{Z} \times \mathbb{Z} \to \mathbb{Z} \) defined by \( (a, b) \mapsto a + b \) is a binary operation on the set of integers \( \mathbb{Z} \).
    \item \textbf{Composition of Symmetries:} Consider a set $X$ equipped with some structure. Now, a symmetry of $X$ is a bijective function from $X$ to itself that preserves the structure of $X$. The operation of composition of symmetries \( \circ : \text{Sym}(X) \times \text{Sym}(X) \to \text{Sym}(X) \) defined by \( (f, g) \mapsto f \circ g \) is a binary operation on the set of symmetries \( \text{Sym}(X) \).
    \item \textbf{Union and Intersection:} Let $S$ be a non-empty set. Then, $\union{}$ and $\intersect{}$ are binary operations on the power set of $S$, denoted by $\mathcal{P}(S)$. We define the operations as follows:
    \[\union{} : \mathcal{P}(S) \times \mathcal{P}(S) \to \mathcal{P}(S)\] defined by \( (A, B) \mapsto A \cup B \)
    \[\intersect{} : \mathcal{P}(S) \times \mathcal{P}(S) \to \mathcal{P}(S)\] defined by \( (A, B) \mapsto A \cap B \)
    \item \textbf{Matrix Multiplication:}For $n \in \N, n\geq 1$, let $GL_n(\R)$ be the set of all invertible $n \times n$ matrices with real entries. Then, the operation of matrix multiplication \( \cdot : GL_n(\R) \times GL_n(\R) \to GL_n(\R) \) defined by \( (A, B) \mapsto A \cdot B \) is a binary operation on the set of invertible matrices \( GL_n(\R) \).
\end{enumerate}
\begin{exercise}
    Is the matrix multiplication operation associative and commutative on the set \( GL_n(\R) \)?
\end{exercise}
\underline{Soln:} The operation of matrix multiplication is associative but not commutative. That is, for any three matrices \( A, B, C \in GL_n(\R) \), we have \( (A \cdot B) \cdot C = A \cdot (B \cdot C) \), but in general, \( A \cdot B \neq B \cdot A \). \\
\begin{exercise}
    which of the following are binary operations on $\N$?
    \begin{enumerate}
        \item \( (a, b) \mapsto a - b \)
        \item \( (a, b) \mapsto \pm \sqrt{ab} \)
        \item \( (a, b) \mapsto 5 \)
    \end{enumerate}
\end{exercise}
\underline{Soln:} Let us analyze each operation:
\begin{enumerate}
    \item The operation \( (a, b) \mapsto a - b \) is not a binary operation on \( \mathbb{N} \) because the result of the operation may not be a natural number. For example, if \( a = 3 \) and \( b = 5 \), then \( a - b = -2 \), which is not in \( \mathbb{N} \).
    \item The operation \( (a, b) \mapsto \pm \sqrt{ab} \) is not a binary operation on \( \mathbb{N} \) because the result of the operation may not be a natural number. For example, if \( a = 2 \) and \( b = 8 \), then \( \sqrt{ab} = \sqrt{16} = 4 \), which is in \( \mathbb{N} \), but the negative root \( -4 \) is not in \( \mathbb{N} \).
    \item The operation \( (a, b) \mapsto 5 \) is a binary operation on \( \mathbb{N} \) because for any natural numbers \( a, b \), the result of the operation is always 5, which is in \( \mathbb{N} \).
\end{enumerate}
\begin{exercise}
    Is the operation on $\R$ defined by \( (a, b) \mapsto \frac{a-b}{2} \) commutative and associative?
\end{exercise} \\
\underline{Soln:} Let us analyze the operation \( (a, b) \mapsto \frac{a-b}{2} \):
\begin{itemize}
    \item \textbf{Commutativity:} The operation is not commutative because in general, \( \frac{a-b}{2} \neq \frac{b-a}{2} \). For example, if \( a = 3 \) and \( b = 5 \), then \( \frac{3-5}{2} = -1 \) and \( \frac{5-3}{2} = 1 \), which are not equal.
    \item \textbf{Associativity:} The operation is not associative because in general, \( \frac{\left(\frac{a-b}{2}\right)-c}{2} \neq \frac{a-\left(\frac{b-c}{2}\right)}{2} \). For example, if \( a = 4, b = 2, c = 0 \), then:
    \[
    \frac{\left(\frac{4-2}{2}\right)-0}{2} = \frac{1-0}{2} = 0.5
    \]
    and
    \[
    \frac{4-\left(\frac{2-0}{2}\right)}{2} = \frac{4-1}{2} = 1.5
    \]
    which are not equal.
\end{itemize}
\subsection{Cayley Tables}
Let S be a non-empty set and $*$ be a binary operation on S. We can represent the operation $*$ using a table called the Cayley table. The rows and columns of the table are labeled with the elements of S, and the entry in the row corresponding to element $a$ and the column corresponding to element $b$ is the result of the operation $a * b$. \\
Let us consider the following example - \\
Let S = $\{ \pm 1, \pm i \}$, where $i$ is the imaginary unit. We define a binary operation $*$ on S as follows -
\begin{center}
\begin{tabular}{c|c|c|c|c|}
    $*$ & $1$ & $-1$ & $i$ & $-i$ \\
    \hline
    $1$ & $1$ & $-1$ & $i$ & $-i$ \\
    $-1$ & $-1$ & $1$ & $-i$ & $i$ \\
    $i$ & $i$ & $-i$ & $-1$ & $1$ \\
    $-i$ & $-i$ & $i$ & $1$ & $-1$ \\
\end{tabular}
\end{center}
\begin{mynote}
    We note the following properties of the above Cayley table -
    \begin{enumerate}
        \item We have 1 as the identity element, since $1 * a = a * 1 = a$ for all $a \in S$.
        \item Every row and every colum contain the identity element, $1 \cdot 1 = 1, -1 \cdot -1 = 1, i \cdot -i = 1, -i \cdot i = 1$
        \item Asssociativity holds, i.e., $(a * b) * c = a * (b * c)$ for all $a, b, c \in S$.
        \item The table is symmetric about the main diagonal, i.e., $a * b = b * a$ for all $a, b \in S$. Thus, the operation is commutative.
        \item No duplicate elements appear in any row or column. If there existed duplicates, then there would be elements without inverses.
    \end{enumerate}
    Based on the above properties, we can conclude that the set S with the operation $*$ forms an abelian group.
\end{mynote}
We now consider a few exercises based on Cayley tables. \\
\begin{exercise}
    Consider the group $(\{e, a, b\}, *)$. Find the cayley table for this group, given that $e$ is the identity element, and that $a * a = b, a * b = e, b * a = e, b * b = a$. Show that the group is abelian.
\end{exercise} \\
\underline{Soln:} We can construct the Cayley table as follows -
\begin{center}
\begin{tabular}{c|c|c|c|}
    $*$ & $e$ & $a$ & $b$ \\
    \hline
    $e$ & $e$ & $a$ & $b$ \\
    $a$ & $a$ & $b$ & $e$ \\
    $b$ & $b$ & $e$ & $a$ \\
\end{tabular}
\end{center}
To show that the group is abelian, we need to verify that the operation is commutative, i.e., \( a * b = b * a \) for all \( a, b \in \{e, a, b\} \). We can see that the cayley table is symmetric about the main diagonal, which implies that the operation is commutative. Therefore, the group is abelian. \\ \\
\begin{exercise}
    Using Cayley tables, find all the groups of order 4 (i.e., groups with 4 elements). Hint: There are 4 possible tables, 3 are identical to each other, and 2 distinct groups exist.
\end{exercise}
\underline{Soln:} Let us denote the elements of the group as \( \{e, a, b, c\} \), where \( e \) is the identity element. We can construct the Cayley tables for the possible groups of order 4. \\
\textbf{Group 1:}
\begin{center}
\begin{tabular}{c|c|c|c|c|}
    $*$ & $e$ & $a$ & $b$ & $c$ \\
    \hline
    $e$ & $e$ & $a$ & $b$ & $c$ \\
    $a$ & $a$ & $b$ & $c$ & $e$ \\
    $b$ & $b$ & $c$ & $e$ & $a$ \\
    $c$ & $c$ & $e$ & $a$ & $b$ \\
\end{tabular}
\end{center}
\textbf{Group 2:}
\begin{center}
\begin{tabular}{c|c|c|c|c|}
    $*$ & $e$ & $a$ & $b$ & $c$ \\
    \hline
    $e$ & $e$ & $a$ & $b$ & $c$ \\
    $a$ & $a$ & $e$ & $c$ & $b$ \\
    $b$ & $b$ & $c$ & $e$ & $a$ \\
    $c$ & $c$ & $b$ & $a$ & $e$ \\
\end{tabular}
\end{center}
These are the two distinct groups of order 4. The other two possible tables are identical to one of these two groups. Thus, we have found all the groups of order 4 using Cayley tables. \\ \\
\begin{exercise}
    Let $(G, *)$ be a group. Prove the following -
    \begin{enumerate}
        \item The identity element is unique.
        \item For each element $a \in G, \exists \text{ unique } b \in G$ such that $a * b = b * a = e$, where $e$ is the identity element.
        \item For all $a, b, c \in G$, if $a * b = a * c$, then $b = c$ (Left Cancellation Law). Similarly, if $b * a = c * a$, then $b = c$ (Right Cancellation Law).
        \item For all $a, b \in G$, the equation $a * x = b$ has a unique solution for $x \in G$. Similarly, the equation $y * a = b$ has a unique solution for $y \in G$.
        \item $(a^{-1})^{-1} = a$ for all $a \in G$ and ($a * b)^{-1} = b^{-1} * a^{-1}$ for all $a, b \in G$.
        \item $\forall a \in G, m, n \in \Z, a^m * a^n = a^{m+n}$, where $a^m$ denotes the $m$-th power of $a$ with respect to the operation $*$. And, $(a^m)^n = a^{mn}$. 
    \end{enumerate}
\end{exercise}
\underline{Soln:} Let us prove each of the statements one by one -
\begin{enumerate}
    \item Suppose there are more than 1 distinct identity elements in the group, say \( e_1 \) and \( e_2 \). By the definition of identity element, we have \( e_1 * a = a \) and \( e_2 * a = a \) for all \( a \in G \). In particular, for \( a = e_2 \), we have \( e_1 * e_2 = e_2 \). Similarly, for \( a = e_1 \), we have \( e_2 * e_1 = e_1 \). Thus, we have \( e_1 = e_2 \), which contradicts our assumption that they are distinct. Therefore, the identity element is unique.
    \item Let \( a \in G \) be an arbitrary element. By the definition of inverse element, there exists an element \( b \in G \) such that \( a * b = e \) and \( b * a = e \), where \( e \) is the identity element. Suppose there exists another element \( c \in G \) such that \( a * c = e \) and \( c * a = e \). Then, we have:
    \[
    b = b * e = b * (a * c) = (b * a) * c = e * c = c
    \]
    Thus, \( b = c \), which shows that the inverse element is unique.
    \item Let \( a, b, c \in G \) such that \( a * b = a * c \). We want to show that \( b = c \). By the definition of inverse element, there exists an element \( a^{-1} \in G \) such that \( a^{-1} * a = e \), where \( e \) is the identity element. Now, we have:
    \[
    a^{-1} * (a * b) = a^{-1} * (a * c) \implies (a^{-1} * a) * b = (a^{-1} * a) * c \implies e * b = e * c \implies b = c  
    \]
    The same logic can be applied to prove the right cancellation law.
    \item We have \( a, b \in G \). We want to show that the equation \( a * x = b \) has a unique solution for \( x \in G \). By the definition of inverse element, there exists an element \( a^{-1} \in G \) such that \( a^{-1} * a = e \), where \( e \) is the identity element. Now, we can multiply both sides of the equation \( a * x = b \) by \( a^{-1} \) from the left:
    \[
    a^{-1} * (a * x) = a^{-1} * b \implies (a^{-1} * a) * x = a^{-1} * b \implies e * x = a^{-1} * b \implies x = a^{-1} * b
    \]
    Thus, the equation \( a * x = b \) has a unique solution for \( x \in G \). Similarly, we can show that the equation \( y * a = b \) has a unique solution for \( y \in G \).
    \item We know that every $a \in G$ has a unique inverse element $a^{-1} \in G$ such that $a * a^{-1} = a^{-1} * a = e$, where $e$ is the identity element. From this, we can deduce that $(a^{-1})^{-1} = a$ for all $a \in G$. \\
    Now, let $a, b \in G$. We want to show that $(a * b)^{-1} = b^{-1} * a^{-1}$. By the definition of inverse element, we have:
    \[(a * b) * (b^{-1} * a^{-1}) = a * (b * b^{-1}) * a^{-1} = a * e * a^{-1} = a * a^{-1} = e\]
    Similarly, we can show that \((b^{-1} * a^{-1}) * (a * b) = e\). Thus, we have shown that \((a * b)^{-1} = b^{-1} * a^{-1}\) for all \(a, b \in G\).
    \item We prove the properties using mathematical induction. \\
    \textbf{Base Case:} For \( n = 1 \), we have \( a^m * a^1 = a^{m+1} \) and \( (a^m)^1 = a^{m \cdot 1} = a^m \), which holds true. \\
    \textbf{Inductive Step:} Assume that the properties hold for some \( n = k \), i.e., \( a^m * a^k = a^{m+k} \) and \( (a^m)^k = a^{m \cdot k} \). We need to show that the properties hold for \( n = k + 1 \). \\
    We have:
    \[a^m * a^{k+1} = a^m * (a^k * a) = (a^m * a^k) * a = a^{m+k} * a = a^{m+k+1}\]
    and
    \[(a^m)^{k+1} = (a^m)^k * a^m = a^{m \cdot k} * a^m = a^{m \cdot k + m} = a^{m \cdot (k + 1)}\]
    Thus, by the principle of mathematical induction, the properties hold for all \( n \in \mathbb{Z} \).
    
\end{enumerate}
We now give some examples of infinite non-abelian groups. \\
\begin{example}
    $GL_n(\R)$, the group of all invertible $n \times n$ matrices with real entries under the operation of matrix multiplication, is an infinite non-abelian group for \( n \geq 2 \). The identity element in this group is the identity matrix, and the inverse of a matrix is its inverse matrix. The operation of matrix multiplication is associative but not commutative, which makes \( GL_n(\R) \) a non-abelian group. \\
    Similarly $SL_n(\R)$, the group of all $n \times n$ matrices with real entries and determinant equal to 1 under the operation of matrix multiplication, is also an infinite non-abelian group for \( n \geq 2 \). The identity element in this group is also the identity matrix, and the inverse of a matrix is its inverse matrix. The operation of matrix multiplication is associative but not commutative, which makes \( SL_n(\R) \) a non-abelian group. \\
    In general, for any infinite field \( F \) and any \( n \geq 2 \), the groups \( GL_n(F) \) and \( SL_n(F) \) are infinite non-abelian groups under the operation of matrix multiplication.
\end{example} 
\begin{example}
    $(\mathbb{Z}/n\mathbb{Z}, +)$, the group of integers modulo \( n \) under the operation of addition, is a finite abelian group. The elements of this group are the equivalence classes of integers modulo \( n \), denoted by \( [0], [1], [2], \ldots, [n-1] \). The identity element in this group is \( [0] \), and the inverse of an element \( [a] \) is \( [-a] \). The operation of addition is associative and commutative, which makes \( \mathbb{Z}/n\mathbb{Z} \) an abelian group. \\
\end{example}
\begin{mynote}
    Each element of $\Z/n\Z$ is an equivalence class of integers. For example, in $\Z/5\Z$, we have the following equivalence classes -
\begin{center}
\begin{align}
    [0] &= \{\ldots -10, -5, 0, 5, 10, \ldots\} \\
    [1] &= \{\ldots -9, -4, 1, 6, 11, \ldots\} \\
    [2] &= \{\ldots -8, -3, 2, 7, 12, \ldots\} \\
    [3] &= \{\ldots -7, -2, 3, 8, 13, \ldots\} \\
    [4] &= \{\ldots -6, -1, 4, 9, 14, \ldots\}
\end{align}
\end{center}
Therefore, each element of $\Z/n\Z$ is itself a countably infinite subset of integers. To make sense of the operation of addition on $\Z/n\Z$, we define the addition of two equivalence classes as follows -
\[[a] + [b] = [a + b]\]
\end{mynote}


Now, another interesting fact is that $\{\Z/n\Z\} \backslash {[0]} = \{[1], [2], \ldots, [n-1]\}$ forms a group under the operation of multiplication for the case of $n=5$. The identity element in this group is $[1]$, and the inverse of an element $[a]$ is $[a^{-1}]$, where $a^{-1}$ is the multiplicative inverse of $a$ modulo $n$. The operation of multiplication is associative and commutative, which makes $\{\Z/n\Z\} \backslash {[0]}$ an abelian group. \\
\begin{mynote}
    Here, we define multiplication as follows - 
    \[[a] \cdot [b] = [a \cdot b]\]
    where $a \cdot b$ is the product of $a$ and $b$ modulo $n$. Note that the operation of multiplication is only defined for the non-zero elements of $\Z/n\Z$, since $[0]$ does not have a multiplicative inverse. Therefore, we consider the set $\{\Z/n\Z\} \backslash {[0]}$ when defining the group under multiplication.
\end{mynote}
We also note, that the same is not a group if $n$ is not prime. 
We state the following result and prove it - 
\begin{mytheorem}
    The group $\{\Z/n\Z\} \backslash {[0]}$ under multiplication is a group if and only if $n$ is prime.
\end{mytheorem}
\begin{proof}
    $(\rightarrow)$Suppose $n$ is not prime and that $\{\Z/n\Z\} \backslash {[0]}$ under multiplication is a group. Since $n$ is not prime, there exist integers $a$ and $b$ such that $1 < a < n$, $1 < b < n$ (as every non prime posseses a unique prime factorization.), and $n = a \cdot b$. Consider the elements $[a]$ and $[b]$ in $\{\Z/n\Z\} \backslash {[0]}$. We have:
    \[[a] \cdot [b] = [a \cdot b] = [n] = [0]\]
    This contradicts the fact that $\{\Z/n\Z\} \backslash {[0]}$ is a group under multiplication, since the product of two non-zero elements cannot be zero in a group. Therefore, if $\{\Z/n\Z\} \backslash {[0]}$ is a group under multiplication, then $n$ must be prime. \\
    $(\leftarrow)$ Suppose $n$ is prime. We need to show that $\{\Z/n\Z\} \backslash {[0]}$ under multiplication is a group. The identity element in this group is $[1]$, since $[1] \cdot [a] = [a] \cdot [1] = [a]$ for all $[a] \in \{\Z/n\Z\} \backslash {[0]}$. Now, let $[a] \in \{\Z/n\Z\} \backslash {[0]}$. Since $n$ is prime, there exists an integer $a^{-1}$ such that $a \cdot a^{-1} \equiv 1 \mod n$ (we prove this seperately). Therefore, we have:
\[[a] \cdot [a^{-1}] = [a \cdot a^{-1}] = [1]\]
Thus, every element in $\{\Z/n\Z\} \backslash {[0]}$ has a multiplicative inverse. The operation of multiplication is associative and commutative, which makes $\{\Z/n\Z\} \backslash {[0]}$ an abelian group under multiplication. Therefore, if $n$ is prime, then $\{\Z/n\Z\} \backslash {[0]}$ under multiplication is a group.
\end{proof}
\begin{mynote}
    To show that if $n$ is prime, there exists an integer $a^{-1}$ such that $a \cdot a^{-1} \equiv 1 \mod n$, consider $a \in \Z, 1 < a < n$. Since $n$ is prime, the greatest common divisor of $a$ and $n$ is 1. We claim there exists integers $x$ and $y$ such that $ax + ny = 1$. This can be shown using the Euclidean algorithm as follows -
\begin{center}
\begin{align*}
    n &= q_1 a + r_1 \\
    a &= q_2 r_1 + r_2 \\
    r_1 &= q_3 r_2 + r_3 \\
    &\vdots \\
    r_{k-2} &= q_k r_{k-1} + r_k \\
    r_{k-1} &= q_{k+1} r_k + 0
\end{align*}
\end{center}
where $q_i$ are the quotients and $r_i$ are the remainders. Since the greatest common divisor of $a$ and $n$ is 1, we have $r_k = 1$. Now, we can express 1 as a linear combination of $a$ and $n$ by back-substituting the equations above:
\begin{center}
\begin{align*}
    1 &= r_k \\
    &= r_{k-2} - q_k r_{k-1} \\
    &= r_{k-2} - q_k (r_{k-3} - q_{k-1} r_{k-2}) \\
    &= (1 + q_k q_{k-1}) r_{k-2} - q_k r_{k-3} \\
    &\vdots \\
    &= x a + y n
\end{align*}
\end{center}
Thus, we have shown that there exist integers $x$ and $y$ such that $ax + ny = 1$. Taking this equation modulo $n$, we get $ax \equiv 1 \mod n$. Therefore, we can take $a^{-1} = x$, which shows that there exists an integer $a^{-1}$ such that $a \cdot a^{-1} \equiv 1 \mod n$.
\end{mynote}
\begin{example}
    Let $A \subseteq \R$ be a non-empty set. The set $C(A)$ of all continous functions from $A$ to $\R$ forms a group under the operation of function addition, defined by $(f + g)(x) = f(x) + g(x)$ for all $f, g \in C(A)$ and $x \in A$. The identity element in this group is the zero function, defined by $0(x) = 0$ for all $x \in A$. The inverse of a function $f \in C(A)$ is the function $-f$, defined by $(-f)(x) = -f(x)$ for all $x \in A$. The operation of function addition is associative and commutative, which makes $C(A)$ an abelian group under function addition.
\end{example}

\subsection{Subgroups}
\begin{mydefinition}{Subgroup}
    Let $(G, *)$ be a group. A subset $H$ of $G$ is called a subgroup of $G$ if $H$ itself forms a group under the operation $*$ inherited from $G$. In other words, $(H, *)$ is a subgroup of $(G, *)$ if the following conditions hold:
    \begin{itemize}
        \item \textbf{Closure:} For all $a, b \in H$, the result of the operation $a * b$ is also in $H$.
        \item \textbf{Identity Element:} The identity element of $G$ is also in $H$.
        \item \textbf{Inverse Element:} For each element $a \in H$, its inverse with respect to the operation $*$ in $G$ is also in $H$.
    \end{itemize}
\end{mydefinition}
We now discuss finite subgroups of $(C^x, \cdot)$, the group of non-zero complex numbers under multiplication. Let $n$ be a positive integer. Consider the set $H_n = \{ e^{2\pi i k/n} : k = 0, 1, 2, \ldots, n-1 \}$. We can verify that $H_n$ is a subgroup of $(C^x, \cdot)$ as follows:
\begin{itemize}
    \item \textbf{Closure:} For any two elements $e^{2\pi i k/n}$ and $e^{2\pi i m/n}$ in $H_n$, their product is given by:
    \[
    e^{2\pi i k/n} \cdot e^{2\pi i m/n} = e^{2\pi i (k+m)/n}
    \]
    Since $k+m$ is an integer, we can write it as $k+m = qn + r$ for some integer $q$ and $0 \leq r < n$. Thus, we have:
    \[
    e^{2\pi i (k+m)/n} = e^{2\pi i (qn + r)/n} = e^{2\pi i q} \cdot e^{2\pi i r/n} = e^{2\pi i r/n}
    \]
    Since $0 \leq r < n$, we have $e^{2\pi i r/n} \in H_n$. Therefore, the product of any two elements in $H_n$ is also in $H_n$, which shows that $H_n$ is closed under multiplication.
    \item \textbf{Identity Element:} The identity element in $(C^x, \cdot)$ is the complex number 1, which can be expressed as $e^{2\pi i \cdot 0/n}$. Since $0$ is an integer, we have $e^{2\pi i \cdot 0/n} = 1 \in H_n$. Thus, the identity element of $(C^x, \cdot)$ is also in $H_n$.
    \item \textbf{Inverse Element:} For each element $e^{2\pi i k/n}$ in $H_n$, its inverse with respect to multiplication is given by:
    \[
    (e^{2\pi i k/n})^{-1} = e^{-2\pi i k/n}
    \]
    We can express this as:
    \[
    e^{-2\pi i k/n} = e^{2\pi i (n-k)/n}
    \]
    Since $n-k$ is an integer and $0 < n-k < n$, we have $e^{2\pi i (n-k)/n} \in H_n$. Therefore, the inverse of each element in $H_n$ is also in $H_n$. Thus, $H_n$ satisfies the conditions for being a subgroup of $(C^x, \cdot)$.
\end{itemize}



\section{Cyclic Groups}
\begin{mydefinition}{Cyclic Group}
    A group \( (G, *) \) is called a cyclic group if there exists an element \( g \in G \) such that every element of \( G \) can be expressed as a power of \( g \) (with respect to the operation \( * \)). In other words, \( G = \{ g^n : n \in \mathbb{Z} \} \), where \( g^n \) denotes the \( n \)-th power of \( g \) with respect to the operation \( * \). The element \( g \) is called a generator of the cyclic group \( G \). 
\end{mydefinition}
We can have both finite and infinite cyclic groups under certain conditions. Consider the following cases - 
\begin{enumerate}
    \item \emphasise{Powers of g are distinct:} If the powers of \( g \) are distinct, i.e., \( g^m \neq g^n \) for all \( m, n \in \mathbb{Z} \) with \( m \neq n \), then the cyclic group \( G = \{ g^n : n \in \mathbb{Z} \} \) is infinite. In this case, the elements of \( G \) are all distinct and there are countably infinite elements in \( G \) (since there is a bijection between \( G \) and \( \mathbb{Z} \) given by \( g^n \leftrightarrow n \)).
    \item \emphasise{Powers of g are not distinct:} If there exist integers \( m, n \in \mathbb{Z} \) with \( m \neq n \) such that \( g^m = g^n \), then the cyclic group \( G = \{ g^n : n \in \mathbb{Z} \} \) is finite. In this case, there exists a positive integer \( k \) such that \( g^k = e \), where \( e \) is the identity element of the group. The order of the cyclic group \( G \) is equal to the smallest positive integer \( k \) such that \( g^k = e \). Note that this set has a bijection with the set of integers modulo \( k \), denoted by \( \mathbb{Z}/k\mathbb{Z} \), which is a finite cyclic group of order \( k \).
    \item \emphasise{Trivial Case:} If \( g = e \), where \( e \) is the identity element of the group, then the cyclic group \( G = \{ g^n : n \in \mathbb{Z} \} \) is trivial and contains only the identity element. In this case, \( G = \{ e \} \) and it is a finite cyclic group of order 1.
\end{enumerate}
\begin{mynote}
    We denote a cyclic group of order \( n \) by \( C_n \) and a cyclic group of infinite order by \( C_\infty \). Thus, we have the following classifications of cyclic groups:
\begin{itemize}
    \item If \( G \) is a cyclic group of order \( n \), then \( G \cong C_n \cong \mathbb{Z}/n\mathbb{Z} \).
    \item If \( G \) is a cyclic group of infinite order, then \( G \cong C_\infty \cong \mathbb{Z} \). 
    \item If \( G \) is a trivial cyclic group, then \( G \cong C_1 \cong \{ e \} \).
\end{itemize}
\end{mynote}
We now define the order of an element in a group.
\begin{mydefinition}{Order of an Element}
    Let \( (G, *) \) be a group and let \( g \in G \) be an element of the group. The order of the element \( g \), denoted by \( \text{ord}(g) \), is defined as the smallest positive integer \( k \) such that \( g^k = e \), where \( e \) is the identity element of the group. If no such positive integer exists, we say that \( g \) has infinite order.
\end{mydefinition}
\begin{example}
    Consider the group \( (\mathbb{Z}/6\mathbb{Z}, +) \), which is a cyclic group of order 6. The elements of this group are \( [0], [1], [2], [3], [4], [5] \). The identity element is \( [0] \). The order of the element \( [1] \) is 6, since \( [1]^6 = [6] = [0] \) and there is no smaller positive integer \( k \) such that \( [1]^k = [0] \). Similarly, the order of the element \( [2] \) is 3, since \( [2]^3 = [6] = [0] \) and there is no smaller positive integer \( k \) such that \( [2]^k = [0] \). The order of the element \( [3] \) is 2, since \( [3]^2 = [6] = [0] \) and there is no smaller positive integer \( k \) such that \( [3]^k = [0] \). The order of the element \( [4] \) is 3, since \( [4]^3 = [12] = [0] \) and there is no smaller positive integer \( k \) such that \( [4]^k = [0] \). The order of the element \( [5] \) is 6, since \( [5]^6 = [30] = [0] \) and there is no smaller positive integer \( k \) such that \( [5]^k = [0] \).
\end{example}
\begin{example}
    Consider the group \( (\mathbb{Z}, +) \), which is an infinite cyclic group. The elements of this group are all integers. The identity element is \( 0 \). The order of any non-zero element \( n \in \mathbb{Z} \) is infinite, since there is no positive integer \( k \) such that \( n^k = 0 \). The order of the element \( 0 \) is 1, since \( 0^1 = 0 \) and there is no smaller positive integer \( k \) such that \( 0^k = 0 \).
\end{example}


\section{Dihedral Groups}
The symmetries of any geometric object form a group under the operation of composition of symmetries. For example, the symmetries of a regular polygon with \( n \) sides form a group called the dihedral group of order \( 2n \), denoted by \( D_n \). The elements of \( D_n \) consist of \( n \) rotations and \( n \) reflections. The group operation is the composition of symmetries, which is associative. The identity element is the identity transformation (which leaves the polygon unchanged), and each element has an inverse (the symmetry that undoes it). The dihedral group \( D_n \) is non-abelian for \( n \geq 3 \), since the composition of a rotation and a reflection does not commute. Let us consider an $n$ sided polygon, and denote the rotation by $r$ and the reflection by $f$. We want to identify the elements of $D_n$ and the relations between them. We have the following relations -
\begin{itemize}
    \item $r^n = e$, since applying the rotation $n$ times results in the identity transformation.
    \item $f^2 = e$, since applying the reflection twice results in the identity transformation.
    \item $f r f^{-1} = r^{-1}$, since conjugating a rotation by a reflection results in the inverse of the rotation.
\end{itemize}
From the above relations, we can deduce that the elements of $D_n$ can be expressed in the form $r^k$ for $k = 0, 1, \ldots, n-1$ (representing the rotations) and $f r^k$ for $k = 0, 1, \ldots, n-1$ (representing the reflections). Thus, the dihedral group $D_n$ has $2n$ elements, and it is generated by the rotation $r$ and the reflection $f$. 
A formal definition of the dihedral group $D_n$ can be given as follows -
\begin{mydefinition}{Dihedral Group}
    The dihedral group of order \( 2n \), denoted by \( D_n \), is the group defined by the presentation:
    \[
    D_n = \langle r, f \mid r^n = e, f^2 = e, f r f^{-1} = r^{-1} \rangle
    \]
    where \( r \) represents a rotation and \( f \) represents a reflection. The elements of \( D_n \) can be expressed as \( r^k \) for \( k = 0, 1, \ldots, n-1 \) (representing the rotations) and \( f r^k \) for \( k = 0, 1, \ldots, n-1 \) (representing the reflections). The group operation is the composition of symmetries.
\end{mydefinition}


\section{Symmetric Groups}
\begin{mydefinition}{Permutation}
    A permutation of a set \( S \) is a bijection from \( S \) to itself. In other words, a permutation is a rearrangement of the elements of \( S \). For example, if \( S = \{1, 2, 3\} \), then one possible permutation is the function that maps \( 1 \) to \( 2 \), \( 2 \) to \( 3 \), and \( 3 \) to \( 1 \). This can be represented in cycle notation as \( (1\ 2\ 3) \).
\end{mydefinition}
We note that the set of all permutations of a set \( S \) forms a group under the operation of composition of functions. The identity element in this group is the identity permutation (which maps each element to itself), and each element has an inverse (the permutation that undoes it). The group of all permutations of a set \( S \) is called the symmetric group on \( S \). If \( S \) has \( n \) elements, then the symmetric group on \( S \) is denoted by \( S_n \) and has \( n! \) elements, since there are \( n! \) ways to rearrange \( n \) distinct objects.
\begin{mydefinition}{Symmetric Group}
    The symmetric group on \( n \) elements, denoted by \( S_n \), is the group of all permutations of \( n \) distinct objects. The elements of \( S_n \) are the bijections from the set \( \{1, 2, \ldots, n\} \) to itself, and the group operation is the composition of these bijections, defined as folows:
    $$\star : S_n \times S_n \to S_n$$
    $$(\sigma, \tau) \mapsto \sigma \cdot \tau$$
    where \( \sigma \cdot \tau \) is the permutation obtained by first applying \( \tau \) and then applying \( \sigma \). The identity element in \( S_n \) is the identity permutation, which maps each element to itself, and the inverse of a permutation is the permutation that undoes it.
\end{mydefinition}
We denote each element of \( S_n \) as follows - 
$${1\ 2\ \ldots\ n} \choose {\sigma(1)\ \sigma(2)\ \ldots\ \sigma(n)}$$
where \( \sigma \) is a permutation of the set \( \{1, 2, \ldots, n\} \). For example, the element of \( S_3 \) that maps \( 1 \) to \( 2 \), \( 2 \) to \( 3 \), and \( 3 \) to \( 1 \) can be represented as:
$${1\ 2\ 3} \choose {2\ 3\ 1}$$

We now define a cycle in a symmetric group.
\begin{mydefinition}{Cycle}
    A $k$-cycle in a symmetric group \( S_n \) is a permutation that maps a subset of \( k \leq n \) distinct elements to each other in a cyclic manner, while leaving the remaining elements fixed. For example, the 3-cycle \( (1\ 2\ 3) \) in \( S_3 \) maps \( 1 \) to \( 2 \), \( 2 \) to \( 3 \), and \( 3 \) to \( 1 \), while leaving any other elements (if they exist) fixed.
\end{mydefinition}
A cycle can be represented in cycle notation, where the elements that are permuted are enclosed in parentheses. For example, the 3-cycle \( (1\ 2\ 3) \) can be represented as:
$${1\ 2\ 3} \choose {2\ 3\ 1}$$
\begin{mydefinition}{Disjoint Cycles}
    Two cycles in a symmetric group \( S_n \) are said to be disjoint if the sets they permute are disjoint. For example, the cycles \( (1\ 2) \) and \( (3\ 4) \) in \( S_4 \) are disjoint, since they permute different sets of elements. On the other hand, the cycles \( (1\ 2) \) and \( (2\ 3) \) in \( S_3 \) are not disjoint, since they both permute the element \( 2 \).
\end{mydefinition}
We now state the following results - 
\begin{mytheorem}[Disjoint Cycles Commute]
    If \( c_1 \) and \( c_2 \) are disjoint cycles in a symmetric group \( S_n \), then they commute with each other. In other words, if \( c_1 \) and \( c_2 \) are disjoint cycles, then \( c_1 c_2 = c_2 c_1 \).
\end{mytheorem}
\begin{proof}
    Let \( c_1 \) and \( c_2 \) be disjoint cycles in \( S_n \). Since \( c_1 \) and \( c_2 \) permute different sets of elements, we can write \( c_1 = (a_1\ a_2\ \ldots\ a_k) \) and \( c_2 = (b_1\ b_2\ \ldots\ b_m) \), where \( a_i \) and \( b_j \) are distinct elements of the set \( \{1, 2, \ldots, n\} \). 

    Now, we want to show that \( c_1 c_2 = c_2 c_1 \). Let \( x \in \{1, 2, \ldots, n\} \). We have three cases to consider:
    \begin{itemize}
        \item If \( x = a_i \) for some \( i \), then \( c_1(x) = a_{i+1} \) (where we take indices modulo \( k \)) and \( c_2(x) = x \). Therefore, we have:
        \[
        (c_1 c_2)(x) = c_1(c_2(x)) = c_1(x) = a_{i+1}
        \]
        and
        \[
        (c_2 c_1)(x) = c_2(c_1(x)) = c_2(a_{i+1}) = a_{i+1}
        \]
        Thus, \( (c_1 c_2)(x) = (c_2 c_1)(x) \).
        
        \item If \( x = b_j \) for some \( j \), then \( c_1(x) = x \) and \( c_2(x) = b_{j+1} \) (where we take indices modulo \( m \)). Therefore, we have:
        \[
        (c_1 c_2)(x) = c_1(c_2(x)) = c_1(b_{j+1}) = b_{j+1}
        \]
        and
        \[
        (c_2 c_1)(x) = c_2(c_1(x)) = c_2(x) = b_{j+1}
        \]
        Thus, \( (c_1 c_2)(x) = (c_2 c_1)(x) \).
        \item If \( x \) is not permuted by either \( c_1 \) or \( c_2 \), then \( c_1(x) = x \) and \( c_2(x) = x \). Therefore, we have:
        \[
        (c_1 c_2)(x) = c_1(c_2(x)) = c_1(x) = x
        \]
        and
        \[
        (c_2 c_1)(x) = c_2(c_1(x)) = c_2(x) = x
        \]
        Thus, \( (c_1 c_2)(x) = (c_2 c_1)(x) \).
    \end{itemize}
    Since \( (c_1 c_2)(x) = (c_2 c_1)(x) \) for all \( x \in \{1, 2, \ldots, n\} \), we conclude that \( c_1 c_2 = c_2 c_1 \). Therefore, disjoint cycles commute with each other.
\end{proof}
\begin{mytheorem}[Decomposition into Disjoint Cycles]
    Every permutation in a symmetric group \( S_n \) can be uniquely expressed as a product of disjoint cycles. In other words, for any permutation \( \sigma \in S_n \), there exist disjoint cycles \( c_1, c_2, \ldots, c_k \) such that \( \sigma = c_1 c_2 \cdots c_k \), and this decomposition is unique up to the order of the cycles.
\end{mytheorem}
\begin{proof}
    Let \( \sigma \in S_n \) be a permutation. We will construct the disjoint cycles that represent \( \sigma \). Start with an element \( i_1 \) in the set \( \{1, 2, \ldots, n\} \). Apply \( \sigma \) to \( i_1 \) to get \( i_2 = \sigma(i_1) \). Continue applying \( \sigma \) to the resulting elements until we return to \( i_1 \). This process will give us a cycle that permutes the elements \( i_1, i_2, \ldots, i_k \), where \( k \leq n\) is the length of the cycle. 

    Next, choose an element \( j_1 \) that has not been permuted by the first cycle. Repeat the same process to find another cycle that permutes a different set of elements. Continue this process until all elements have been permuted by some cycle. Since each element can only be permuted by one cycle, the cycles we obtain are disjoint.
    Now, we show we can express \( \sigma \) as the product of these disjoint cycles. Let \( c_1, c_2, \ldots, c_m \) be the disjoint cycles obtained from the process above. Since the cycles are disjoint, they commute with each other. Therefore, we can write:
    \[\sigma = c_1 c_2 \cdots c_m\]
\end{proof}
\begin{mytheorem}[Order of a $k$-cycle]
    The order of a $k$-cycle in a symmetric group \( S_n \) is \( k \). In other words, if \( \sigma \) is a $k$-cycle, then \( \sigma^k = e \), where \( e \) is the identity permutation, and there is no smaller positive integer \( m < k \) such that \( \sigma^m = e \).
\end{mytheorem}
\begin{proof}
    Let \( \sigma = (a_1\ a_2\ \ldots\ a_k) \) be a $k$-cycle in \( S_n \). We want to show that \( \sigma^k = e \) and that there is no smaller positive integer \( m < k \) such that \( \sigma^m = e \). 

    First, we compute \( \sigma^2 \):
    \[
    \sigma^2(a_1) = \sigma(a_2) = a_3, \quad \sigma^2(a_2) = \sigma(a_3) = a_4, \quad \ldots, \quad \sigma^2(a_{k-1}) = \sigma(a_k) = a_1, \quad \sigma^2(a_k) = \sigma(a_1) = a_2
    \]
    Thus, \( \sigma^2 = (a_1\ a_3\ a_5\ \ldots)(a_2\ a_4\ a_6\ \ldots) \). Continuing this process, we find that:
    \[
    \sigma^3(a_1) = a_4, \quad \sigma^3(a_2) = a_5, \quad \ldots
    \]
    After applying \( k-1 \) times, we have:
    \[
    \sigma^{k-1}(a_1) = a_k, \quad \sigma^{k-1}(a_2) = a_1, \quad \ldots
    \]
    Finally, we compute \( \sigma^k(a_i) = a_i \) for all \( i = 1, 2, ..., k \), which shows that \( \sigma^k = e \).

    Now, suppose there exists a positive integer \( m < k \) such that \( \sigma^m = e \). This would imply that applying the cycle \( m \) times results in the identity permutation. However, since the cycle permutes \( k \) distinct elements in a cyclic manner, it would take exactly \( k \) applications of the cycle to return to the original configuration. Therefore, there cannot be any smaller positive integer \( m < k \) such that \( \sigma^m = e\).
\end{proof}
\begin{mytheorem}[Order of a Product of Disjoint Cycles]
    The order of a product of disjoint cycles in a symmetric group \( S_n \) is the least common multiple (LCM) of the lengths of the cycles. In other words, if \( \sigma = c_1 c_2 \cdots c_m \) is a product of disjoint cycles, where \( c_i \) is a cycle of length \( k_i \), then the order of \( \sigma \) is given by:
    \[
    \text{ord}(\sigma) = \text{lcm}(k_1, k_2, \ldots, k_m)
    \]
\end{mytheorem}
\begin{proof}
    Let \( \sigma = c_1 c_2 \cdots c_m \) be a product of disjoint cycles, where \( c_i \) is a cycle of length \( k_i \). Since the cycles are disjoint, they commute with each other. Therefore, we can write:
    \[
    \sigma^n = (c_1 c_2 \cdots c_m)^n = c_1^n c_2^n \cdots c_m^n
    \]
    We want to find the smallest positive integer \( n \) such that \( \sigma^n = e \). This is equivalent to finding the smallest positive integer \( n \) such that \( c_i^n = e \) for all \( i = 1, 2, ..., m \).

    Since \( c_i \) is a cycle of length \( k_i \), we know from the previous theorem that \( c_i^{k_i} = e \) and there is no smaller positive integer \( m < k_i \) such that \( c_i^m = e \). Therefore, for each cycle \( c_i \), the order is \( k_i \).

    To ensure that \( c_i^n = e \) for all \( i = 1, 2, ..., m \), we need to find the smallest positive integer \( n \) that is a multiple of each \( k_i \). This is precisely the definition of the least common multiple (LCM). Thus, we have:
    \[
    n = \text{lcm}(k_1, k_2, ..., k_m)
    \]
    Therefore, the order of the product of disjoint cycles \( \sigma \) is given by:
    \[
    \text{ord}(\sigma) = n = \text{lcm}(k_1, k_2, ..., k_m)
    \]
\end{proof}
\begin{mytheorem}[Number of $k$-cycles in $S_n$]
    The number of $k$-cycles in the symmetric group \( S_n \) is given by the formula:
    \[(k-1)! \cdot \binom{n}{k}\]
    where \( \binom{n}{k} \) is the binomial coefficient representing the number of ways to choose \( k \) elements from a set of \( n \) elements, and \( (k-1)! \) accounts for the number of distinct cycles that can be formed with those \( k \) elements.
\end{mytheorem}
\begin{proof}
    Choose \( k \) distinct elements from the set \( \{1, 2, \ldots, n\} \) to form the cycle. The number of ways to choose \( k \) elements from \( n \) is given by the binomial coefficient:
    \[
    \binom{n}{k}
    \]
    Once we have chosen the \( k \) elements, we can arrange them in a cycle. The number of ways to arrange \( k \) distinct elements in a cycle is given by \( (k-1)! \), since we can fix one element and permute the remaining \( k-1 \) elements in any order. Therefore, the total number of \( k \)-cycles in \( S_n \) is given by:
    \[
    (k-1)! \cdot \binom{n}{k}
    \]
\end{proof}
We now define a two cycle in a symmetric group.
\begin{mydefinition}{Two Cycle}
    A two cycle (also known as a transposition) in a symmetric group \( S_n \) is a cycle that permutes exactly two elements while leaving the remaining elements fixed. For example, the two cycle \( (1\ 2) \) in \( S_3 \) maps \( 1 \) to \( 2 \), \( 2 \) to \( 1 \), and leaves \( 3 \) fixed.
\end{mydefinition}
\begin{mynote}
    We note that for all transpositions $\sigma$, we have $\sigma^2 = e$, since applying the transposition twice results in the identity permutation. It immediately follows that $\sigma^{-1} = \sigma$
\end{mynote}
\begin{mydefinition}{Parity of a Permutation}
    The parity of a permutation \( \sigma \in S_n \) is defined as follows:
    \begin{itemize}
        \item If \( \sigma \) can be expressed as a product of an even number of transpositions, then \( \sigma \) is called an even permutation, and its parity is defined to be 0.
        \item If \( \sigma \) can be expressed as a product of an odd number of transpositions, then \( \sigma \) is called an odd permutation, and its parity is defined to be 1.
    \end{itemize}
\end{mydefinition}
\begin{mytheorem}[Any cycle can be expressed as a product of transpositions]
    Every cycle in a symmetric group \( S_n \) can be expressed as a product of transpositions. In particular, a $k$-cycle can be expressed as a product of \( k-1 \) transpositions.
\end{mytheorem}
\begin{proof}
    Let \( \sigma = (a_1\ a_2\ \ldots\ a_k) \) be a $k$-cycle in \( S_n \). We can express \( \sigma \) as a product of transpositions as follows:
    \[
    \sigma = (a_1\ a_k)(a_1\ a_{k-1}) \cdots (a_1\ a_2)
    \]
    To verify that this product of transpositions is equal to the original cycle \( \sigma \), we can apply both sides to an arbitrary element \( x \in \{1, 2, \ldots, n\} \) and check that they yield the same result. We have three cases to consider:
    \begin{itemize}
        \item If \( x = a_1 \), then the right-hand side maps \( a_1 \) to \( a_2 \), which is the same as the left-hand side.
        \item If \( x = a_i \) for some \( i = 2, 3, \ldots, k \), then the right-hand side maps \( a_i \) to \( a_{i+1} \) (where we take indices modulo \( k \)), which is the same as the left-hand side.
        \item If \( x \) is not one of the elements in the cycle, then both sides map \( x \) to itself.
        \item Thus, we have shown that \( \sigma = (a_1\ a_k)(a_1\ a_{k-1}) \cdots (a_1\ a_2) \), which is a product of \( k-1 \) transpositions. Therefore, every cycle can be expressed as a product of transpositions.
    \end{itemize}
\end{proof}
\begin{corollary}
    Every permutation in a symmetric group \( S_n \) can be expressed as a product of transpositions. In particular, if a permutation \( \sigma \) can be expressed as a product of disjoint cycles, then we can express each cycle as a product of transpositions, and thus express \( \sigma \) as a product of transpositions.
\end{corollary}
\begin{proof}
    Let \( \sigma \in S_n \) be a permutation. From the theorem on decomposition into disjoint cycles, we can express \( \sigma \) as a product of disjoint cycles:
    \[
    \sigma = c_1 c_2 \cdots c_m
    \]
    where each \( c_i \) is a cycle. From the previous theorem, we know that each cycle \( c_i \) can be expressed as a product of transpositions. Therefore, we can write:
    \[
    \sigma = (t_{i1} t_{i2} \cdots t_{i_{l_i}})_{i=1}^m
    \]
    where each \( t_{ij} \) is a transposition. This shows that \( \sigma \) can be expressed as a product of transpositions.
\end{proof}
\begin{mytheorem}[Number of odd and even permutations in $S_n$]
    The number of even permutations in the symmetric group \( S_n \) is equal to the number of odd permutations, and both are equal to \( \frac{n!}{2} \). In other words, there are \( \frac{n!}{2} \) even permutations and \( \frac{n!}{2} \) odd permutations in \( S_n \).
\end{mytheorem}
\begin{proof}
    We can establish a bijection between the set of even permutations and the set of odd permutations in \( S_n \) by fixing a transposition \( t \) and defining a function \( f: S_{n_e} \to S_{n_0} \) as follows:
    \[
    f(\sigma) = t \cdot \sigma
    \]
    for any even permutation \( \sigma \in S_{n_e} \). 

    We need to show that \( f \) is a bijection. First, we show that \( f \) is injective. Suppose \( f(\sigma_1) = f(\sigma_2) \) for some \( \sigma_1, \sigma_2 \in S_n \). This means that:
    \[
    t \cdot \sigma_1 = t \cdot \sigma_2
    \]
    Multiplying both sides by \( t^{-1} = t \), we get:
    \[
    t^{-1} \cdot (t \cdot \sigma_1) = t^{-1} \cdot (t \cdot \sigma_2)
    \]
    which simplifies to:
    \[
    \sigma_1 = \sigma_2
    \]
    Thus, \( f \) is injective.

    Next, we show that \( f \) is surjective. Let \( y \in S_{n_o} \) be any permutation. We want to find a permutation \( x \in S_{n_e} \) such that \( f(x) = y \). We can choose \( x = t^{-1} y = t y\), since \( t^{-1} = t\). Then we have:
    \[
    f(x) = f(t y) = t (t y) = (t t) y = e y = y
    \]
    Thus, for every permutation \( y \in S_{n_o}\), there exists a permutation \( x = t y\) such that \( f(x) = y\). Therefore, \( f\) is surjective.

    Since \( f\) is both injective and surjective, it is a bijection. This means that the number of even permutations in \( S_n\) is equal to the number of odd permutations in \( S_n\). Since the total number of permutations in \( S_n\) is \( n!\), we have:
    \[\text{Number of even permutations} + \text{Number of odd permutations} = n!\]
    Let \( E \) be the number of even permutations and \( O \) be the number of odd permutations. We have:
    \[E = O = \frac{n!}{2}\]
    Hence proved. 
\end{proof}
\begin{mytheorem}[$S_{n_e}$ is a subgroup of $S_n$]
    The set of even permutations in the symmetric group \( S_n \), denoted by \( S_{n_e} \), forms a subgroup of \( S_n \). In other words, \( S_{n_e} \) is closed under the group operation (composition of permutations), contains the identity element, and contains the inverse of each of its elements.
\end{mytheorem}
\begin{proof}
    To show that \( S_{n_e} \) is a subgroup of \( S_n \), we need to verify the following three properties:
    \begin{itemize}
        \item \textbf{Closure:} If \( \sigma_1, \sigma_2 \in S_{n_e} \), then \( \sigma_1 \cdot \sigma_2 \in S_{n_e} \).
        \item \textbf{Identity:} The identity permutation \( e \) is in \( S_{n_e} \).
        \item \textbf{Inverses:} If \( \sigma \in S_{n_e} \), then the inverse permutation \( \sigma^{-1} \) is also in \( S_{n_e} \).
    \end{itemize}

    \textbf{Closure:} Let \( \sigma_1, \sigma_2 \in S_{n_e} \) be two even permutations. By definition, \( \sigma_1 \) can be expressed as a product of an even number of transpositions, and \( \sigma_2 \) can also be expressed as a product of an even number of transpositions. When we compose \( \sigma_1 \) and \( \sigma_2 \), the resulting permutation \( \sigma_1 \cdot \sigma_2 \) can be expressed as a product of the transpositions from \( \sigma_1 \) and the transpositions from \( \sigma_2 \). Since both \( \sigma_1 \) and \( \sigma_2 \) are even, the total number of transpositions in the composition \( \sigma_1 \cdot \sigma_2 \) is also even. Therefore, \( \sigma_1 \cdot \sigma_2 \in S_{n_e} \).
    \textbf{Identity:} The identity permutation \( e \) can be expressed as a product of zero transpositions, which is an even number. Therefore, \( e \in S_{n_e} \).
    \textbf{Inverses:} Let \( \sigma \in S_{n_e} \) be an even permutation. Since \( \sigma \) can be expressed as a product of an even number of transpositions, we can write:
    \[\sigma = t_1 t_2 \cdots t_{2k}\]
    where \( t_i \) are transpositions and \( 2k \) is an even number. The inverse of \( \sigma \) can be expressed as:
\[\sigma^{-1} = t_{2k}^{-1} t_{2k-1}^{-1} \cdots t_1^{-1}\]
    Since each transposition is its own inverse (i.e., \( t_i^{-1} = t_i \)), we have:
    \[\sigma^{-1} = t_{2k} t_{2k-1} \cdots t_1\]
    This is also a product of an even number of transpositions, so \( \sigma^{-1} \in S_{n_e} \).

    Since \( S_{n_e} \) satisfies closure, contains the identity element, and contains the inverses of its elements, we conclude that \( S_{n_e} \) is a subgroup of \( S_n \).
\end{proof}
\begin{mytheorem}[Transpositions as generators of $S_n$]
    The symmetric group \( S_n \) is generated by the set of transpositions given by $(1 \ 2), (2 \ 3), \ldots, (n-1 \ n)$. In other words, every permutation in \( S_n \) can be expressed as a product of these transpositions.
\end{mytheorem}
\begin{proof}
    We will show that any transposition can be expressed as a product of the transpositions \( (1 \ 2), (2 \ 3), \ldots, (n-1 \ n) \). Since every permutation can be expressed as a product of transpositions, it will follow that every permutation can be expressed as a product of the given transpositions.

    Let \( (i \ j) \) be a transposition where \( 1 \leq i < j \leq n \). We want to express \( (i \ j) \) as a product of the transpositions \( (1 \ 2), (2 \ 3), \ldots, (n-1 \ n) \).

    We can achieve this by first moving \( i \) to position \( 1 \) and \( j \) to position \( 2 \) using the given transpositions, then applying the transposition \( (1 \ 2) \), and finally moving the elements back to their original positions. The sequence of transpositions can be written as follows:
    \[
    (i\ j) = (1\ i)(1\ j)(1\ i)
    \]
    This expression shows that we can express any transposition \( (i\ j) \) as a product of the transpositions \( (1\ 2), (2\ 3), \ldots, (n-1\ n) \).

    Since every permutation in \( S_n \) can be expressed as a product of transpositions, and we have shown that every transposition can be expressed as a product of the given transpositions, it follows that every permutation in \( S_n \) can be expressed as a product of the transpositions \( (1\ 2), (2\ 3), \ldots, (n-1\ n) \). Therefore, these transpositions generate the symmetric group \( S_n \).
\end{proof}
\begin{mynote}
    We note that a set of $n-2$ or less transpositions cannot generate $S_n$. For example, the set of transpositions $\{(1\ 2), (2\ 3), \ldots, (n-2\ n-1)\}$ cannot generate $S_n$. \\
    We also note that there are $n^{n-2}$ different ways to generate $S_n$ using $n-1$ transpositions.
\end{mynote}

\subsection{Cycle Type}
\begin{mydefinition}{Cycle Type}
    The cycle type of a permutation \( \sigma \in S_n \) is a way to classify permutations based on the lengths of the cycles in their disjoint cycle decomposition. The cycle type is often represented as a partition of \( n \) that indicates how many cycles of each length are present in the decomposition. For example, if a permutation has two 1-cycles (fixed points), one 2-cycle, and one 3-cycle, its cycle type can be represented as \( (1^2, 2^1, 3^1) \).
\end{mydefinition}
We note the following facts:
\begin{itemize}
    \item $\sum_{k=1}^{n} k \cdot m_k = n$, where \( m_k \) is the number of cycles of length \( k \) in the cycle decomposition of the permutation.
    \item $k_1 + k_2 + \ldots + k_m$ forms a partition of \( n \), where \( k_i \) is the length of the \( i \)-th cycle arranged in ascending order in the cycle decomposition of the permutation.
\end{itemize}
\begin{mytheorem}[Number of permutations with a cycle type]
    The number of permutations in \( S_n \) with a given cycle type \( (1^{m_1}, 2^{m_2}, \ldots, n^{m_n}) \) is given by the formula:
    \[
    \frac{n!}{\prod_{k=1}^{n} (k^{m_k} \cdot m_k!)}
    \]
    where \( m_k \) is the number of cycles of length \( k \) in the cycle decomposition of the permutation.
\end{mytheorem}
\begin{proof}
    The number of possible arrangements of n symbols is given by $n!$. However, we need to account for the indistinguishability of cycles of the same length and the indistinguishability of elements within each cycle. For each cycle of length \( k \), there are \( k \) ways to arrange the elements within the cycle, and for \( m_k \) cycles of length \( k \), there are \( m_k! \) ways to arrange the cycles themselves. Therefore, we need to divide \( n! \) by the product of \( k^{m_k} \) (to account for the arrangements within each cycle) and \( m_k! \) (to account for the arrangements of the cycles). We need to consider the overcounting for all $k$ from 1 to \( n \), which gives us the formula:
    \[
    \frac{n!}{\prod_{k=1}^{n} (k^{m_k} \cdot m_k!)}
    \]
    Hence proved.
\end{proof}


\end{document}



